\clearpage
\section{Koremenos, Lipson, and Snidal -- The Rational Design of International Institutions}

\subsection{Research question and rational-design approach}

\emph{why are international institutions designed in such different ways?} Their central presumption is that states, together with other relevant international actors, construct institutions in order to advance their goals and to address the particular cooperation problems they face $\Rightarrow$ Institutional design $\Rightarrow$ outcome of purposive and strategic choice.

The argument is located in the rational-choice tradition. Earlier rationalist work concentrated mainly on whether cooperation could emerge under anarchy. Repeated interaction and the "shadow of the future" showed how self-interested states might sustain cooperation without centralized enforcement.

They define international institutions broadly as explicit arrangements negotiated among international actors that prescribe, prohibit, or authorize behavior.

Tacit bargains are excluded because the focus is on deliberately negotiated institutional arrangements.

Nonstate actors such as firms, NGOs, and international organizations may also participate in institutional design.

The rational-design approach differs from both realist and constructivist simplifications. Against a strong realist view, institutions are not merely empty reflections of power: states invest heavily in institutional bargaining precisely because rules affect outcomes and often persist even when the distribution of power changes. Against a strong constructivist interpretation, institutions should not be treated simply as autonomous producers of norms: they are also consciously created arrangements designed by actors with strategic purposes. The authors do not claim that rational design explains every feature of every institution, but they argue that it provides a systematic framework for identifying recurring relationships between cooperation problems and institutional forms.

\subsection{From cooperation theory to institutional design}

The authors begin from the insight that repeated interaction makes cooperation possible but emphasize that the real world contains complications that simple repeated Prisoner's Dilemma models tend to suppress.

The authors explain design variation through a set of independent variables:

\begin{enumerate}
    \item \textbf{Distribution problems} arise because there are usually multiple possible cooperative outcomes and actors disagree about how the gains from cooperation should be divided. Cooperation is therefore not simply a choice between cooperation and defection; it often requires bargaining over which cooperative equilibrium will be selected.

    \item \textbf{Enforcement problems} arise because actors may have \textbf{incentives to defect}.

    \item The \textbf{number of actors} matters. Large groups complicate bargaining, monitoring, communication, and collective action. They also raise problems of asymmetry when participants differ significantly in power, wealth, or capabilities.

    \item \textbf{Uncertainty} IOs can help manage these different forms of uncertainty by collecting information, monitoring behavior, providing procedures for adjustment, and reducing transaction costs.
\end{enumerate}

The central move of the article is to connect these cooperation problems to observable institutional design choices. Their rules must be broadly \emph{incentive compatible}: actors must have reasons to create, maintain, and comply with the institution over time. It means that, in the long run, participation must remain preferable to abandonment for the relevant actors.

Institutional development is also path dependent: previous rules constrain later options. Rational design can nevertheless operate within such historical constraints through conscious modification, imitation of successful arrangements, or selection among existing institutions.

\subsection{Five dimensions of institutional design}

The framework identifies five main dependent variables: \textbf{membership}, \textbf{scope}, \textbf{centralization}, \textbf{control}, and \textbf{flexibility}. These dimensions do not exhaust all institutional variation, but they are substantively important, measurable, and applicable across very different institutional settings.

\begin{enumerate}
    \item \textbf{Membership.} Membership concerns who belongs to an institution and under what conditions. Institutions may be universal, regional, highly restricted, or open to nonstate actors. Membership determines who receives the benefits of cooperation, who bears costs, and which actors participate in enforcement. NATO enlargement illustrates how membership can have both practical and symbolic consequences. Restricting membership can improve compliance by excluding actors whose incentives or capacities are inconsistent with the agreement, but inclusion can be useful when excluded states could undermine the arrangement or when broad participation is necessary to solve a distribution problem.
    \[enforcement\ problems\ \uparrow\ \Rightarrow\ membership \downarrow\]
    \[distribution\ problems\ \uparrow\ \Rightarrow\ membership \uparrow\]

    \item \textbf{Scope: range of issues included in an agreement}. Issue linkage can facilitate agreement by creating possibilities for side payments and package deals. The Nuclear Nonproliferation Treaty, for instance, linked non-acquisition of nuclear weapons to access to peaceful nuclear technology.

    \textbf{Consequence:} issue linkage becomes more attractive when distributional conflict is severe because a broader agenda permits side payments and cross-issue trades.
    Scope can also expand when enforcement across issues becomes mutually reinforcing.

    \item \textbf{Centralization.} Centralization refers to whether important tasks are assigned to a focal institutional body. The authors use the concept more broadly than centralized enforcement. Centralization does not necessarily imply a surrender of sovereignty: institutions can combine centralized information or adjudication with decentralized enforcement by member states.

    \textbf{Consequence:} When monitoring is difficult, a central body can gather information and distinguish noncompliance from noise.

    When many actors are involved, centralization can reduce communication and transaction costs, coordinate behavior, and organize collective punishment.

    \item \textbf{Control:} how collective decisions are made and who has influence over them (voting rules). An institution can be highly centralized in information collection while leaving control equally distributed among members, or it can remain relatively decentralized while giving certain members disproportionate voting power.

    \textbf{Consequence:} asymmetries in power and uncertainty about distributional consequences affect the allocation of decision-making authority. Powerful states may demand disproportionate control when they make larger contributions or bear greater risks. At the same time, states may seek vetoes or other safeguards when they are uncertain about how future decisions will affect them. Institutional control rules are therefore closely linked to bargaining power and anticipated distributional consequences.

    \item \textbf{Flexibility.} Flexibility concerns the capacity of an institution to accommodate new circumstances. The authors distinguish \emph{adaptive} flexibility from \emph{transformative} flexibility. Adaptive devices, such as escape clauses, allow temporary deviations while preserving the broader arrangement. Transformative devices allow more fundamental change. Flexibility is especially important when future conditions are uncertain, but it can create new costs by encouraging opportunistic use of exceptions or repeated bargaining.

    \textbf{Consequence:}
    \begin{itemize}
        \item [(a)] flexibility $\uparrow$ when uncertainty about the state of the world $\uparrow$. When future conditions are difficult to predict, rigid commitments are risky.
        \item [(b)] flexibility $\uparrow$ when distributional problems $\uparrow$. because flexibility can reduce the stakes of the initial bargain.
        \item [(c)] flexibility $\downarrow$ when number of actors $\uparrow$. Renegotiation becomes more costly as the number of participants rises.
    \end{itemize}
\end{enumerate}

\subsection{Design interactions and institutional complexity}

Institutional features are \textbf{substitutes} when different arrangements solve the same problem: for instance, an escape clause provides decentralized flexibility, whereas a centralized body empowered to authorize exceptions can solve the same problem through centralized discretion.

Institutional features are \textbf{complements} when one feature strengthens another. Restrictive membership may support enforcement, while centralized monitoring makes it easier to determine whether members are actually complying.
Finally, design principles can \textbf{conflict}. Severe enforcement problems may favor restricted membership, while severe distribution problems favor broader membership.

Interactions also occur across institutions. Institutions may be nested within one another, or one arrangement may be selected because it provides more suitable enforcement mechanisms than another. The shift of intellectual-property issues toward the WTO rather than the World Intellectual Property Organization illustrates institutional selection based on design advantages.

\subsection{Contribution and limitations}

The overall conclusion is that international institutions should be understood as negotiated responses to strategic problems (or independent variables). Institutional design matters because different rules change the incentives, information, and bargaining possibilities available to states. The rational-design project thus reconnects the detailed study of institutional arrangements with the theoretical study of international cooperation.

\clearpage
\section{Leeds -- Do Alliances Deter Aggression? The Influence of Military Alliances on the Initiation of Militarized Interstate Disputes}

\subsection{The puzzle: do alliances produce peace or war?}

Leeds addresses a long-standing dispute over whether military alliances deter aggression or instead make war more likely. Previous large-$N$ studies produced inconsistent findings, leading some scholars to conclude that alliances have little independent effect on conflict. Leeds argues that this conclusion follows from treating all alliances as if they were equivalent. In reality, alliance treaties contain different promises that send different signals to potential challengers and targets. Some commitments deter aggression, while others reassure or embolden states considering the use of force. Aggregating all alliances into a single category can therefore hide their effects.

The central theoretical claim is informational. Formal alliance treaties provide state leaders with credible information about how outside states are likely to behave if a crisis occurs. Alliances function as costly signals because states incur costs in negotiating them and may face future costs if they fail to honor their obligations. Existing evidence that alliance commitments are fulfilled in a large majority of applicable wartime cases makes \textbf{treaty provisions relevant to strategic calculations}.

\textbf{Leeds defines alliances as written agreements between at least two sovereign states that contain promises of defensive assistance, offensive assistance, neutrality, nonaggression, or consultation in circumstances involving military conflict or crisis.} The key point is that the precise \emph{content} and \emph{conditions} of these commitments matter. Treaties often specify particular adversaries, geographical areas, or triggering events. They therefore cannot be treated as unconditional ``blank checks.''

\subsection{Three hypotheses about treaty content}

\begin{enumerate}
    \item \textbf{Defensive commitments to the potential target} $\Rightarrow$ aggression $\downarrow$. If a challenger knows that attacking a target will activate a commitment by an outside ally to provide military assistance, the challenger must anticipate a multilateral rather than bilateral conflict.

    \item \textbf{Offensive commitments to the potential challenger} $\Rightarrow$ aggression $\uparrow$. A state contemplating aggression is more likely to proceed when it expects allied military support. An offensive alliance does not necessarily create revisionist preferences, but it changes the expected costs and probability of success for a challenger that already has reasons to use force. Thus, challengers with allies committed to participate in an offensive operation should be more likely to initiate militarized disputes.

    \item \textbf{Neutrality commitments to the challenger} $\Rightarrow$ aggression $\uparrow$. A neutrality pact can assure the challenger that a third party will not intervene on the target's side. This reduces the risk that the conflict will expand into an unfavorable coalition war. The challenger therefore gains information not about direct military assistance, but about the absence of opposing intervention. \textbf{Leeds predicts that relevant neutrality commitments increase the likelihood of dispute initiation.}
\end{enumerate}

\subsection{Research design}

The empirical analysis uses the Alliance Treaty Obligations and Provisions (ATOP) dataset, which codes the specific content of alliance treaties. The sample covers politically relevant directed dyad-years from 1816 to 1944. A directed dyad distinguishes, for example, Britain initiating against Russia from Russia initiating against Britain, which is necessary because the theory concerns the decision of a specific potential challenger to initiate a dispute against a specific target.

The dependent variable is whether the challenger initiates a Militarized Interstate Dispute (MID) against the target in a given year.

The three main independent variables identify whether:
\begin{itemize}
    \item the target has a relevant defensive ally
    \item the challenger has a relevant offensive ally
    \item the challenger has an ally whose neutrality commitment would prevent intervention on behalf of the target.
\end{itemize}

\textbf{Secret defensive alliances} are \textbf{excluded} from the main specification because an unknown commitment cannot deter an adversary through information transmission.

\subsection{Findings}

\begin{enumerate}
    \item All three variables are statistically significant in the expected directions and add substantial explanatory power to the baseline model.

    \item These results help \textbf{reconcile earlier contradictory findings}. If defensive alliances make conflict less likely while offensive and neutrality commitments make it more likely, an \textbf{aggregate variable for ``alliances'' can easily produce a weak or inconsistent average effect}. The apparent empirical puzzle is therefore partly an operationalization problem.

    \item Leeds concludes that alliances do affect international behavior and can function as instruments of general deterrence, but their consequences depend on treaty design. The correct question is not whether alliances as a category cause peace or war. Instead, scholars must \textbf{identify the information contained in specific agreements and the incentives that those commitments create for particular actors}.

    \item The findings also have policy implications. Defensive guarantees can deter challengers by increasing the expected likelihood of outside intervention. At the same time, alliances can create moral-hazard-like effects by reassuring partners that they will receive assistance or by guaranteeing that potential opponents will remain neutral. Policymakers therefore need to consider not only whom they ally with but precisely what obligations they accept and how those obligations will be interpreted by third parties.

    \item More broadly, the article supports the view that \textbf{formal international agreements matter because their detailed provisions convey credible information}. Treaty content should therefore be treated as a substantive variable rather than as legal detail.
\end{enumerate}

\clearpage
\section{Morrow --- Alliances: Why Write Them Down?}

\subsection{The central question}

States can cooperate militarily without a formal alliance, and an alliance cannot legally force a state to intervene in war. Why, then, do states formalize military relationships in written treaties?

The answer is that written alliances can alter both expectations before war and incentives during war. They may increase the credibility of future intervention through two distinct mechanisms
\begin{itemize}
    \item \textbf{Signaling:} allows other states to infer that allies share sufficiently strong interests to support one another.
    \item \textbf{Commitment:} changes the incentives of alliance members so that intervention becomes more likely than it would have been without the treaty. 
\end{itemize}

These mechanisms explain \textbf{why formalization can matter even in an anarchic system} where agreements ultimately have to be self-enforcing.

\textbf{Alignment:} It exists when \textbf{common interests are sufficiently obvious} that actors do not need to formalize their relationship.

\textbf{Alliance:} It is an \textbf{explicit commitment} specifying expected behavior under defined wartime contingencies.

\textbf{Tradeoff between security and autonomy}: States may gain security through promised support, but they sacrifice some freedom because they accept obligations toward allies.

\subsection{The decision to intervene in war}

To understand what alliances change, Morrow develops a simple expected-utility framework for third-party intervention. Suppose A and B are at war and C prefers an A victory. C compares the expected utility of staying out with the expected utility of intervening on A's side. The intervention decision depends on three main considerations.

The first is C's \textbf{marginal military contribution}: how much does C's participation increase the probability that A defeats B, weighted by how strongly C prefers A's victory to B's? Intervention becomes more attractive when C can substantially affect the outcome or has a strong interest in which side wins.

The second is any \textbf{private gain from participation}. C may receive side payments, territory, influence, or a favorable place in the postwar settlement. Such gains provide motives for intervention beyond the collective security effect of defeating the opponent.

The third is the possibility of \textbf{additional losses if C intervenes and its side loses}. A state that enters a war can suffer territorial, political, or even existential losses beyond the direct costs of fighting.

Because the values associated with these considerations are often private information, other states cannot know with certainty whether a third party will intervene. This is the credibility problem that alliances can address.

\subsection{Alliances as signals}
For an alliance to signal commitment, forming or maintaining it must impose costs that less committed states are unwilling to bear.

One possible cost is \textbf{entrapment}: an alliance may increase the risk that a state is drawn into a war initiated or escalated by its partner. A state with a high aversion to war may therefore avoid a treaty that could expose it to this danger. Signing can consequently reveal a relatively high willingness to fight for the ally.

\textbf{Policy coordination} can also make alliances costly and therefore informative. To preserve credibility, allies often align their foreign policies on issues where their preferences are not identical. This can force each state to abandon positions it would otherwise prefer. The willingness to bear such peacetime costs signals that the relationship rests on substantial shared interests.

\textbf{Military coordination} produces another costly signal. Such specialization creates vulnerabilities if the alliance collapses, so states unwilling to sustain the relationship are less likely to accept these arrangements.

\textbf{Secret provisions} do not directly inform adversaries, but they can communicate limits and commitments between the allies themselves. Offensive alliances may be kept secret because each partner needs reassurance that the other accepts its intended revision of the status quo without alerting the intended target.

\subsection{Alliances as commitment devices}

Commitment differs from signaling because it changes behavior rather than merely revealing preexisting incentives. Morrow emphasizes two mechanisms.

First, \textbf{military coordination} can raise the probability of joint victory. Intervention becomes more valuable because the ally's marginal military contribution is larger. This directly changes the expected-utility calculation at the moment of crisis.

Second, \textbf{audience costs} can make reneging more costly. External audiences may punish states that acquire reputations for abandoning allies, while domestic audiences may punish leaders who violate explicit international commitments.

Audience-cost arguments nevertheless require caution. External reputation may not transfer easily across cases because interests, capabilities, and costs differ across crises. Failure to defend one partner does not automatically imply failure to defend another. Morrow therefore calls for more explicit theories explaining when reputational inferences are valid.

Commitment mechanisms also create the danger of \textbf{entrapment}. If an ally becomes confident that support will be forthcoming, it may take greater risks or initiate disputes. States respond by limiting the conditions under which alliance obligations apply. Yet alliance terms are rarely fully specified. Complete specificity would reduce entrapment but would also tell adversaries exactly which issues fall outside the alliance, weakening deterrence. Ambiguity can therefore be strategically valuable by creating a ``halo'' of possible support beyond the explicitly covered contingencies.

\subsection{Critique of neorealism and balance-of-power theory}

Morrow then evaluates neorealist theories of alliances. Traditional balance-of-power arguments distinguish status quo states from revisionist states and treat security against a revisionist power as a public good. Status quo states may nevertheless fail to balance because they prefer to pass the costs of resistance to others, producing \emph{buckpassing}; alternatively, they may align with the threatening power, producing \emph{bandwagoning}. The security dilemma and the danger of chain-ganging further complicate alliance behavior.

Morrow argues that this framework is incomplete because all states possess both status quo and revisionist interests. A state may defend the existing order on some issues while desiring change on others. The relevant question is therefore not whether a state is inherently status quo or revisionist, but how strongly it values particular changes and whether it is willing to pay the costs of pursuing them.

This point changes the logic of balancing. If states can gain private benefits from intervention, then the provision of security is not a pure public-good problem. Side payments and revisionist opportunities can make states more willing to join coalitions than neorealist theory suggests. Similarly, the central security problem is often not simply balancing capabilities but inferring the scope and intensity of others' intentions. Britain and France in the 1930s did not need to be convinced to resist a Germany bent on unlimited conquest; they needed to determine whether Hitler's demands were limited or expansive.

Morrow therefore argues that balance-of-power theory is more accurately a theory of strategic interests and alignments than a theory of why some relationships are formalized as alliances. Formal alliance theory must explain the informational and commitment functions of written agreements.

\subsection{Arms, burden sharing, alliance management, and asymmetry}

The article surveys several additional alliance problems. One is the \textbf{tradeoff between arms and allies}. A state can improve security by increasing its own military capabilities or by obtaining external commitments. Both are costly in different ways. Arms require material resources, while alliances impose autonomy costs and risks of entrapment or abandonment. In principle, states should substitute between the two at the margin. Empirically, however, large security shocks may cause both armament and alliance formation to increase simultaneously, making the underlying substitution difficult to observe in aggregate data.

A related question is \textbf{burden sharing}. When deterrence is a collective benefit, \textbf{smaller members may free ride} on larger members because they bear the full cost of their own military expenditures but receive only part of the collective deterrent benefit.

Alliance \textbf{management and duration} depend on changing capabilities and interests. When a state's own capabilities increase, the security value of its alliance may decline; when an ally's capabilities decline, the benefits of the alliance may also fall. Regime changes can alter national interests and increase the probability of alliance termination. Snyder's framework adds that alliance management involves balancing fears of abandonment against fears of entrapment.

Morrow also distinguishes \textbf{symmetric and asymmetric alliances}. In asymmetric alliances, a weaker state typically receives security while offering other concessions to a stronger protector, such as political influence, bases, or alignment on foreign-policy issues. Such alliances often last longer because the partners' benefits are complementary and because the weaker state remains dependent on protection. Another interpretation is that the stronger ally demands influence over the weaker ally in order to control the risk of entrapment.

\subsection{Domestic politics and future research}

Domestic politics shapes alliance policy. Empirical work suggests that democratic alliances tend to last longer and that democracies are comparatively reliable partners.

\textbf{The core answer to the question ``why write them down?'' is therefore that formal alliances can do things informal common interests cannot do as effectively.} They can reveal the depth of shared interests, create costs for policy divergence or noncompliance, improve joint military effectiveness, and expose leaders to political penalties for reneging. Alliances matter not because international law mechanically enforces them, but because formalization changes information and incentives.

\clearpage
\section{Lall --- Beyond Institutional Design: Explaining the Performance of International Organizations}

\subsection{The performance puzzle}

Why do some international organizations perform much better than others? Performance is defined as the extent to which an IO achieves its stated objectives while operating cost-effectively and responding to a broad range of public and private stakeholders.

The contrast between the Food and Agriculture Organization (FAO) and the World Food Programme (WFP) is especially striking. Both are nearly universal organizations concerned with global food security, both have governing bodies composed of state representatives, and both have thousands of employees. Yet the WFP is rated among the best-performing IOs while the FAO is rated among the worst.

Instead, the central problem is the tendency of \textbf{states themselves} to interfere with IO policies in order to advance narrow national interests. The principal, rather than the agent, is the main source of moral hazard.

\subsection{Time inconsistency and the principal's moral hazard}

The argument begins from a standard delegation framework. States create IOs because they share interests that can be advanced more effectively through centralized cooperation. After delegation individual states often discover that particular IO policies threaten domestic constituencies, geopolitical objectives, commercial interests, or other national priorities. They then have incentives to pressure the IO to deviate from the broader organizational objectives that motivated delegation in the first place.

This produces a \textbf{time-inconsistency problem}. At time $t$, states prefer an institution capable of pursuing collective goals. At time $t+1$, the same states may prefer to distort its policies for parochial purposes. If states retain effective control over the institution, the result can be gridlock, inefficient resource allocation, politically selective programming, and poor responsiveness to stakeholders.

This logic reverses the conventional public-choice and principal-agent interpretation. In standard ``rogue-agency'' accounts, international bureaucrats exploit information asymmetries and weak monitoring to maximize budgets, influence, or bureaucratic interests at the expense of member states.

The solution is \textbf{policy autonomy}: the capacity of IO officials to determine which mandate-related problems to address and how to address them without state interference. An IO with greater policy autonomy should therefore perform better, not worse.

\subsection{De jure versus de facto policy autonomy}

A crucial distinction is between \textbf{de jure} and \textbf{de facto} policy autonomy. De jure autonomy is the autonomy granted by the formal rules used to design the institution. De facto autonomy is the autonomy the organization actually enjoys in practice.

Policy autonomy has three main components

\begin{enumerate}
    \item IO officials must possess \textbf{agenda-setting power}: the ability to propose policies, draft budgets, and prepare the agendas of governing bodies.
    
    \item Policy autonomy depends on the ability of individual states to \textbf{veto} decisions. Majority voting limits unilateral obstruction, whereas consensus decision making effectively gives every state a veto. 
    
    \item autonomy increases when IOs have \textbf{access to nonstate or independently earned financing}
\end{enumerate}

Lall therefore predicts that de facto policy autonomy will be strongly and positively associated with performance, whereas de jure autonomy will be only weakly associated with both de facto autonomy and performance.

\textbf{Institutional design matters less than the way institutions actually operate.}

\subsection{Sources of de facto autonomy: institutionalized alliances}

If formal design cannot guarantee autonomy, what does? Lall identifies two main sources: \textbf{institutionalized alliances} with actors above and below the state and the \textbf{technical complexity} of IO activities.

Institutionalized alliances are sustained collaborative relationships between IOs and actors such as NGOs, businesses, public-private partnerships, and other international organizations. These partners share important policy objectives with the IO and provide complementary resources, information, legitimacy, and organizational capacity.

Partners can defend IO autonomy in several ways. They can lobby governments domestically, publicize the costs of state interference, mobilize counter-coalitions internationally, provide expertise and logistical support, and contribute financing that reduces dependence on states. These relationships increase the political and material costs governments face when attempting to manipulate IO policy.

Not every partnership has this effect. Lall argues that protective alliances require three features.

\begin{enumerate}
    \item Cooperation must be \textbf{deep} rather than symbolic: partners need meaningful roles in policy formulation, implementation, monitoring, or enforcement.
    \item Partners' preferences must be \textbf{aligned} with the IO's organizational objectives.
    \item Partner capabilities must be \textbf{complementary} and sufficiently large relative to the IO's own resources.
\end{enumerate}

\subsection{Technical complexity as a source of autonomy}

Technical complexity refers to the amount and specificity of specialized training and professional expertise required to formulate policy.

These information asymmetries protect autonomy in several ways:
\begin{itemize}
    \item Governments are more likely to defer to IO officials' agenda-setting authority
    \item less capable of organizing effective opposition
    \item Technical expertise can also generate independent income through the sale of specialized services, publications, advice, or financial products, further reducing state leverage.
\end{itemize}

This argument again reverses the conventional principal-agent perspective. Information asymmetry is normally treated as a source of agency loss because it weakens principals' ability to monitor agents.

\subsection{Measuring IO performance and autonomy}

Lall constructs the first cross-organizational quantitative dataset explicitly focused on IO performance, covering fifty-three organizations. Performance is measured using official evaluations by the British, Australian, Danish, Dutch, and Swedish governments. These assessments combine stakeholder judgments, independent evaluations, audits, quantitative indicators, and IO reporting.

The performance concept has three dimensions: \textbf{goal attainment}, \textbf{cost-effectiveness}, and \textbf{responsiveness to diverse stakeholders}. The author emphasizes four measurement requirements. Measures should be inclusive of different stakeholder perspectives, incorporate objective evidence, cover multiple dimensions of performance, and display empirical coherence across indicators. The government assessments meet these conditions reasonably well, and the different national ratings are strongly correlated, reducing concern that the results reflect the narrow preferences of one donor government.

De facto policy autonomy is measured using information on agenda-setting authority in practice, actual governing-body decision procedures, and actual sources of financing. De jure autonomy uses parallel indicators drawn from formal constitutional and procedural rules. The two are only modestly correlated. In roughly four-fifths of the organizations, formal and actual autonomy differ, demonstrating that institutional design is an imperfect guide to institutional practice.

\subsection{Statistical results}

The quantitative evidence strongly supports the main argument. De facto policy autonomy is positively associated with IO performance across the major government assessments and across the separate dimensions of performance.

De jure policy autonomy, by contrast, performs poorly as an explanation.

Institutionalized alliances and technical complexity are positively and significantly associated with de facto policy autonomy.

By contrast, de jure autonomy remains positive but statistically insignificant in models explaining actual autonomy.

\textbf{the autonomy an institution is designed to possess has surprisingly little bearing on the autonomy it actually enjoys.}

\subsection{FAO and WFP: similar designs, different performance}

The comparative case study of the FAO and WFP illustrates the causal mechanisms. The FAO was designed with substantial formal autonomy. Its director-general was supposed to propose policies, draft the budget, and shape governing-body agendas; most decisions were formally subject to majority voting; and the organization had some capacity to earn revenue independently. In practice, however, member states came to dominate its policy-making process. States used committees, financial pressure, delays in contributions, and exit threats to promote national interests. Formal rules provided little protection against this informal intervention.

The result has been low de facto autonomy and weak performance. The FAO has struggled to pursue politically sensitive reforms and has increasingly concentrated on less controversial activities such as data collection, advice, and voluntary codes. Funding problems and institutional gridlock have limited field operations. Administrative costs have remained high, and the organization has been responsive primarily to governmental rather than diverse stakeholders. The persistence of global undernourishment despite major reductions in extreme poverty is used as evidence of limited success in fulfilling the FAO's core mission.

The WFP presents the opposite pattern. It was also designed with considerable autonomy, but unlike the FAO it has preserved substantial autonomy in practice. WFP officials dominate the preparation of strategic plans, budgets, reports, and project proposals, while states often play a more limited approval role. The WFP has also diversified its financing more successfully and has developed strong relationships with nonstate and international partners.

This autonomy has allowed the WFP to resist donor attempts to politicize aid allocation. During the Cold War, for example, it supported projects in Soviet-aligned countries despite opposition from major Western donors. More generally, WFP aid is strongly related to recipient need and less closely related to donor geopolitical interests than bilateral food aid. The organization has also shifted toward local and regional food procurement and cash or voucher assistance despite resistance from some donor countries that benefited from traditional tied food aid.

The contrast between the two food organizations also supports the proposed sources of autonomy. The WFP has deeper and more capable external partnerships than the FAO and works in a highly operational environment in which specialized expertise and logistics strengthen staff authority. These conditions make interference more difficult and increase access to nonstate resources. The comparison therefore demonstrates how similar formal designs can produce very different actual governance structures and performance outcomes.

\subsection{Theoretical and policy implications}

The article challenges the functionalist assumption that rational institutional design will naturally generate efficient outcomes. Formal rules often fail to determine how institutions operate because sovereign states can circumvent them. Researchers should therefore devote greater attention to \emph{de facto} institutional characteristics rather than relying on constitutional rules as proxies for actual authority.

The argument also revises principal-agent theory in international relations. The principal's moral hazard deserves as much attention as agency slack. International bureaucrats may sometimes resist states not because they are pathological or self-serving but because they are protecting the collective objectives that states themselves established at the moment of delegation.

Finally, Lall connects literatures on IO autonomy, nonstate actors, and epistemic expertise. External partnerships can increase rather than constrain bureaucratic discretion when partner organizations share institutional objectives and help defend the IO against state interference. Technical expertise similarly creates political protection by generating informational advantages.

The policy implication is that strengthening formal safeguards alone is unlikely to improve performance substantially. Effective reform should focus on building the conditions that sustain autonomy in practice: deeper partnerships with capable and aligned nonstate or supranational actors, diversified financing, and the cultivation of professional expertise. The central conclusion is therefore that institutional performance depends less on the rules written into an organization's founding documents than on the political and informational environment that determines whether those rules can actually be respected.
